%% HAND-WRITTEN fragment -- NOT generated by maestro2tex (raw Spectre
%% transients, not a Maestro results database; there is no maestro2tex config
%% for any *_eis_* fragment).  Precedent: fig_eis_error.tex, fig_pixtop_noise.tex.
%%
%% Drop-in source: ~/chips/castalia/simruns/eisfigs2_20260826/build/
%% fig4_settling.tex, standalone preamble and in-figure note removed -- the note
%% is the caption here.
%% Data: data/eis_settle_*.dat (panel a) and data/eis_excess_*.dat (panel b),
%% copied from eisfigs2_20260826/data/fig4a_set_*.dat and fig4c_set_*.dat,
%% reduced there from ~/chips/castalia/simruns/eisdeep_20260826/csv/
%% settle_{set_c1u_f1k,set_c1u_f10k,set_c330n_f10k,set_c100n_f10k}_win{1,8}.csv
%% -- the four compliant record-start sweeps.  data/eis_excess_exp.dat is the
%% exp(-t/tau) guide.  The per-run ac reference and tau are the header line of
%% each source csv; window dependence is settle/window_dependence.txt.
%%
%% Panel (b) is the ABSOLUTE excess over each run's own asymptote, not the
%% normalised form.  eisfigs2's fig4b_*.dat carry the same collapse normalised
%% to each run's excess at t_start = 0, on which all four land on exp(-t/tau)
%% exactly; the absolute axis is plotted instead because it also carries the
%% 0.1 % design rule, which a normalised axis cannot.
%%
%% COLOUR IS ORDINAL, NOT CATEGORICAL.  The series are ordered by T/tau, which
%% is a magnitude, so they take a one-hue blue ramp light -> dark with T/tau.
%% The two runs that share T/tau = 1.010 SHARE A COLOUR and are separated by
%% mark shape only: same colour means same T/tau, and the fact that the two
%% marks lie on top of each other IS the result.  Do not recolour them apart.
%%
%% Key numbers.  Each run asymptotes to the ac analysis of its own load and code
%% within 0.0026-0.0091 percentage points (-0.0076 / -0.0029 / -0.0026 /
%% -0.0091).  The two runs at matched T/tau = 1.010, whose tau and f differ 10x,
%% agree to 0.69 % at every point.  The decay is exp(-t_start/tau): excess
%% 0.582 % at 3 tau, 0.078 % at 5 tau, 0.0105 % at 7 tau -> wait 5 tau for
%% 0.1 %, 7 tau for 0.01 %.  Amplitude is set by the window, not by tau or f: at
%% the onset one cycle leaves -10.6 %, eight cycles -2.3 %, sixteen -1.2 %.
%% Compliance: max |V_OUT - vcm| 0.358-0.447 V of 0.80 V on all four.
%% tau = C_dl*(R_u||R_ct) = 99 us / 0.99 ms / 9.9 ms for fresh / typical /
%% fouled.
%%
%% THE ARCHIVED +16.4 % / +74 % ARE DELIBERATELY NOT PLOTTED, and the caption
%% says so.  They are reproduced exactly by this campaign (+16.355 % /
%% +73.960 % for the 8-cycle window at the onset) and they are SATURATION: gain
%% code N = 20, whose own |Z|min is 1.67x this load, with V_OUT running 1.25 V
%% against the +-0.80 V window.  They asymptote at +8.64 % / +8.08 % against ac
%% answers of 1.08 % / 2.40 %, i.e. they do not converge at all.  Their data is
%% in eisfigs2_20260826/data/fig4_settling.csv flagged saturated = 1.
%%
%% Other caveats: the phase channel carries no useful settling signal (asymptote
%% 0.0071-0.0114 deg from the ac answer, whole excursion 0.018-0.075 deg, the
%% same order), which is why both panels are magnitude only.  The sliding window
%% is exact, not an approximation: two separate simulations that start the
%% excitation 4.2 ms and 12.2 ms into the transient agree with the sliding
%% windows at the same offset to 8e-11 percentage points.  Panel (a)'s x-axis is
%% truncated at 8 tau; the records run to 39 tau and are flat there.  Nominal
%% corner, 37 C, no mismatch, no noise sources -- everything plotted is a
%% deterministic error.
%% Library cell names are kept out of the rendered prose.
\providecommand{\MaestroRoot}{include/analog/}
%% Ordinal ramp -- deliberately NOT the chapter's categorical mtxA..mtxE
%% slots, and named apart from them so that a \providecolor here can never
%% be pre-empted by an earlier fragment's definition of the same name.
%% (mtxD is EDA100 everywhere else in this chapter; defining the dark end of
%% this ramp as mtxD silently rendered two of the four series yellow.)
\providecolor{mtxOL}{HTML}{86B6EF}
\providecolor{mtxOM}{HTML}{3987E5}
\providecolor{mtxOD}{HTML}{184F95}
\begin{figure}[!htb]
  \centering
  {\pgfplotsset{compat=1.18}%
  \pgfplotsset{
    stax/.style={
      width=0.46\linewidth, height=5.4cm,
      grid=both,
      major grid style={line width=0.2pt, draw=black!14},
      minor grid style={line width=0.2pt, draw=black!7},
      axis line style={line width=0.4pt, draw=black!45},
      tick style={line width=0.4pt, draw=black!45},
      tick align=outside,
      label style={font=\small}, tick label style={font=\small},
      title style={font=\small},
      every axis plot/.append style={line join=round, line width=0.9pt,
                                     mark size=1.5pt},
      xmin=-0.3, xmax=8.2, xtick={0,1,2,3,4,5,6,7,8},
      xminorticks=false,
      xlabel={Record start after excitation onset~[\(\tau\)]},
      clip=false,
    }}%
  \begin{tikzpicture}
    \begin{axis}[stax,
      ymin=-11.6, ymax=2.6, ytick={-10,-8,-6,-4,-2,0}, yminorticks=false,
      ylabel={\(|Z|\) magnitude error~[\si{\percent}]},
      title={\textbf{(a) the record converges on the answer}},
    ]
      \addplot[mtxOL, mark=triangle*, mark options={fill=mtxOL}]
        table {\MaestroRoot data/eis_settle_c1u_f10k.dat};
      \addplot[mtxOM, mark=diamond*, mark options={fill=mtxOM}]
        table {\MaestroRoot data/eis_settle_c330n_f10k.dat};
      \addplot[mtxOD, mark=*, mark options={fill=mtxOD}]
        table {\MaestroRoot data/eis_settle_c1u_f1k.dat};
      \addplot[mtxOD, mark=square*, mark options={fill=mtxOD}]
        table {\MaestroRoot data/eis_settle_c100n_f10k.dat};
      \addplot[black!55, line width=0.7pt, densely dotted, no marks,
        domain=-0.3:8.2, samples=2] {0.30};
      \node[font=\scriptsize, anchor=south east] at (axis cs:8.2,0.45)
        {small-signal answer, \SIrange{0.21}{0.34}{\percent}};
    \end{axis}
  \end{tikzpicture}\hfill
  \begin{tikzpicture}
    \begin{axis}[stax, ymode=log,
      ymin=2e-3, ymax=400, ytick={1e-2,1e-1,1e0,1e1},
      yticklabels={0.01,0.1,1,10}, yminorticks=false,
      ylabel={\(|\)transient excess\(|\)~[\si{\percent}]},
      title={\textbf{(b) one curve in \(t/\tau\)}},
      legend style={font=\scriptsize, draw=none, fill=white,
                    fill opacity=0.85, text opacity=1, row sep=0.5pt,
                    at={(0.98,0.98)}, anchor=north east,
                    cells={anchor=west}},
    ]
      \addplot[black!55, densely dashed, line width=0.7pt, no marks,
        forget plot] table {\MaestroRoot data/eis_excess_exp.dat};
      \addplot[mtxOL, mark=triangle*, mark options={fill=mtxOL}]
        table {\MaestroRoot data/eis_excess_c1u_f10k.dat};
      \addlegendentry{\SI{1}{\micro\farad}, \SI{10}{\kilo\hertz},
                      \(T/\tau = 0.10\)}
      \addplot[mtxOM, mark=diamond*, mark options={fill=mtxOM}]
        table {\MaestroRoot data/eis_excess_c330n_f10k.dat};
      \addlegendentry{\SI{330}{\nano\farad}, \SI{10}{\kilo\hertz},
                      \(T/\tau = 0.31\)}
      \addplot[mtxOD, mark=*, mark options={fill=mtxOD}]
        table {\MaestroRoot data/eis_excess_c1u_f1k.dat};
      \addlegendentry{\SI{1}{\micro\farad}, \SI{1}{\kilo\hertz},
                      \(T/\tau = 1.01\)}
      \addplot[mtxOD, mark=square*, mark options={fill=mtxOD}]
        table {\MaestroRoot data/eis_excess_c100n_f10k.dat};
      \addlegendentry{\SI{100}{\nano\farad}, \SI{10}{\kilo\hertz},
                      \(T/\tau = 1.01\)}
      \draw[black, densely dotted, line width=0.6pt]
        (axis cs:-0.3,0.1) -- (axis cs:8.2,0.1);
      \draw[black, densely dotted, line width=0.6pt]
        (axis cs:5,2e-3) -- (axis cs:5,0.1);
      \node[font=\scriptsize, anchor=south west, align=left]
        at (axis cs:5.1,0.115) {\textbf{wait \(5\tau\) for \SI{0.1}{\percent}}};
    \end{axis}
  \end{tikzpicture}}
  \caption[{Impedance error against how late the demodulation window
  starts.}]{Impedance-magnitude error against how late the demodulation window
  starts, in units of the electrode's dc step time constant
  \(\tau = C_{\mathrm{dl}}(R_u \parallel R_{\mathrm{ct}})\), for four conditions
  spanning \(10\times\) in \(\tau\) (\SI{99}{\micro\second} to
  \SI{990}{\micro\second}) and \(10\times\) in excitation frequency, all on the
  correctly gain-ranged code 0 and all in compliance
  (\(\max |V_{\mathrm{OUT}} - v_{\mathrm{cm}}|\)
  \SIrange{0.358}{0.447}{\volt} of \SI{0.80}{\volt}); each run asymptotes to the
  small-signal analysis of its own load and code within \numrange{0.003}{0.009}
  percentage points, which is what validates the demodulator. \textbf{(b)} plots
  the excess over each run's own asymptote, and it is a function of \(t/\tau\)
  alone: colour is ordered by \(T/\tau\) rather than naming a condition, so the
  two runs whose \(\tau\) and frequency differ \(10\times\) at matched
  \(T/\tau\) share one colour, and that their marks lie on top of each
  other --- \SI{0.69}{\percent} apart at every point --- is the result. All four
  decay as \(\exp(-t/\tau)\), giving the rule \textbf{wait \(5\tau\) for
  \SI{0.1}{\percent} and \(7\tau\) for \SI{0.01}{\percent}}; amplitude, not
  decay rate, is set by the demodulation window, so eight cycles taken at the
  onset still leave \SI{-2.3}{\percent} and sixteen \SI{-1.2}{\percent} and
  averaging is not a substitute for waiting. \textbf{The published
  \(+16.4\,\%\) and \(+74\,\%\) are deliberately not plotted}: those records sit
  on a gain code whose own minimum measurable impedance is \(1.67\times\) this
  load, with the output at \SI{1.25}{\volt} against a
  \(\pm\SI{0.80}{\volt}\) window, and they asymptote at \(+8.6\,\%\) and
  \(+8.1\,\%\) against small-signal answers of \(1.08\,\%\) and \(2.40\,\%\) ---
  saturation, not settling. On code 0 the same eight-cycle record at the onset
  reads \SI{-2.00}{\percent} and \SI{-1.32}{\percent}, \(8\times\) and
  \(56\times\) smaller and the opposite sign.}
  \label{fig:eis-settling}
\end{figure}
